%\subsection{\texorpdfstring{\colorbox{yellow}{Time-Averaging Output}}%
%                                             {Time-Averaging Output}}
\subsection{Time-Averaging Output}%
\label{SEC:Files:Time_Averaging_Output}
\key*{}[time averaging]

%%%%%%%%%%%%%%%%%%%%%%%%%%%%%%%%%%%%%%%%%%%%%%%%%%%%%%%%%%%%%%%%%%%%%%%%%%%%%%%%

When performing large-eddy simulations, time-averaged variables are
often needed to aid in post-processing.
%
Storing all of the instantaneous flow variables at every sampling time
would permit one to compute any time-averaged quantity at the end of the
simulation by re-processing the data, but would also require a prohibitive
amount of storage.
%
Instead, these time-averaged variables can be computed on-the-fly
during each iteration of the solution process and output as a
separate set of restart and CGNS files whenever those are written.

%%%%%%%%%%%%%%%%%%%%%%%%%%%%%%%%%%%%%%%%%%%%%%%%%%%%%%%%%%%%%%%%%%%%%%%%%%%%%%%%

{%BEG_newcommand_scope
\newcommand{\ux}{u_x}
\newcommand{\ur}{u_r}
\newcommand{\ut}{u_\theta}
\newcommand{\cp}{c_p}
 
The list of variables that are computed and saved in the
time-averaged restart file is as follows.
%
\begin{align*}
   & \ReyAve{\rho}      && \ReyAve{\rho^2} 
\\ & \ReyAve{T}         && \ReyAve{T^2} && \ReyAve{T_t}
\\ & \ReyAve{p}         && \ReyAve{p^2} && \ReyAve{p_t}
%   & \ReyAve{\rho^2}    && \ReyAve{p^2} && \ReyAve{T^2}
%\\ & \ReyBar{\rho  }    && \ReyBar{p}   && \ReyBar{T}
\\ & \ReyAve{\rho u}    && \ReyBar{u}   
\\ & \ReyAve{\rho v}    && \ReyBar{v}   && \ReyAve{\rho \ur} && \ReyAve{\ur}
\\ & \ReyAve{\rho w}    && \ReyBar{w}   && \ReyAve{\rho \ut} && \ReyAve{\ut}
\\ & \ReyAve{\rho e_t}
  && \ReyAve{\rho T}
\displaybreak[2]
\\ & \ReyAve{\ux \ux}   && \ReyAve{\rho \ux \ux}
  && \ReyAve{\ux  T}    && \ReyAve{\rho \ux T}
\\ & \ReyAve{\ur \ur}   && \ReyAve{\rho \ur \ur}
  && \ReyAve{\ur  T}    && \ReyAve{\rho \ur T}
\\ & \ReyAve{\ut \ut}   && \ReyAve{\rho \ut \ut}
  && \ReyAve{\ut  T}    && \ReyAve{\rho \ut T}
\\ & \ReyAve{\ux \ur}   && \ReyAve{\rho \ux \ur}
\\ & \ReyAve{\ux \ut}   && \ReyAve{\rho \ux \ut}
\\ & \ReyAve{\ur \ut}   && \ReyAve{\rho \ur \ut}
\end{align*}
%
Note that these are the time averages of instantaneous flow variables, and
the higher order terms do not directly represent the turbulent fluctuations.
%
However, this set of terms provides the ability during post-processing to
compute conservative and primitive mean-flow variables,
turbulent stresses, and turbulent heat fluxes
in either the Reynolds-averaged or Favre-averaged form.
%
For example,
\begin{align}
%   Reynolds stress
    - \ReyBar{\rho} \ReyAve{\ReyFlx{u}_x \ReyFlx{u}_\theta}
 =& - \ReyBar{\rho}
      \left(
        \ReyAve{u_x u_\theta}
      - \ReyAve{u}_x \ReyAve{u}_\theta
      \right)
  & 
%   Favre stress
    - \ReyAve{\rho \FavFlx{u}_x \FavFlx{u}_\theta}
 =& -
      \left(
        \ReyAve{\rho u_x u_\theta}
      - \ReyBar{\rho} \FavAve{u}_x \FavAve{u}_\theta
      \right)
  & 
%   Favre velocity
      \FavAve{u}_x
 =&   \ReyAve{\rho u_x} / \ReyBar{\rho}
    \label{EQN:Time_Ave_PostProc_Equations}
 \\
%   Reynolds heat flux
      \ReyBar{\rho} \ReyAve{\ReyFlx{u}_x \ReyFlx{h}}
 =&   \ReyBar{\rho} \cp
      \left(
        \ReyAve{u_x T}
      - \ReyBar{u}_x \ReyBar{T}
      \right)
  &
%   Favre heat flux
      \ReyAve{\rho \FavFlx{u}_x \FavFlx{h}}
 =&   \cp
      \left(
        \ReyAve{\rho u_x T}
      - \ReyBar{\rho} \FavBar{u}_x \FavBar{T}
      \right)
  &
%   Favre velocity
      \FavAve{u}_\theta
 =&   \ReyAve{\rho u_\theta} / \ReyBar{\rho}
    \nonumber
 \\
 &&
 &&
%   Favre temperature
      \FavBar{T}
 =&   \ReyAve{\rho T} / \ReyBar{\rho}
    \nonumber
\end{align}
%
If the input
\inp[nmlTimeAve::TimeAveVelIsAxisymm]{Time\_Ave\_Vel\_Is\_Axisymm}
is false, then the variables $(\ux,\ur,\ut)$ in the statistics above
are replaced with $(u,v,w)$ and the Cartesian stresses and
heat fluxes are retained instead.
%
If a rotating reference frame is used, then the output will also
be with respect to that reference frame.

%%%%%%%%%%%%%%%%%%%%%%%%%%%%%%%%%%%%%%%%%%%%%%%%%%%%%%%%%%%%%%%%%%%%%%%%%%%%%%%%

Even though all of the previously mentioned variables are computed and saved in
the time-averaged restart file, most users are primarily interested in those
needed to compute the Reynolds-averaged quantities.
%
Therefore only the subset of variables in
\tref{TAB:CGNS_Time_Ave_Output_Variables} is output to the time-averaged CGNS
file.
%
Changing the output list requires modifying subroutine {\tt
get\_cgns\_time\_ave\_variables} and recompiling the code.
%
In the CGNS time-averaged solution files, there are several variables whose
names consist of only two letters, \forexample \texttt{TT}, \texttt{uv},
\texttt{rf}.
%
These two letter combinations represent the fluctuating Reynolds-averaged terms,
with the previously mentioned examples representing
$\left(\texttt{TT},\texttt{uv},\texttt{rf}\right)=
\left(\ReyAve{\ReyFlx{T}\ReyFlx{T}},
      \ReyAve{\ReyFlx{u}\ReyFlx{v}},
      \ReyAve{\ReyFlx{\ur}\ReyFlx{\ut}}\right)$.
 
%%%%%%%%%%%%%%%%%%%%%%%%%%%%%%%%%%%%%%%%%%%%%%%%%%%%%%%%%%%%%%%%%%%%%%%%%%%%%%%%

\begin{table}[h]
\begin{center}
\caption[Time-averaged variables output to CGNS.]%
        {Time-averaged variables output to CGNS based on
         \inp[nmlTimeAve::TimeAveVelIsAxisymm]{Time\_Ave\_Vel\_Is\_Axisymm}.}
\label{TAB:CGNS_Time_Ave_Output_Variables}
\begingroup % limit scope of local commands
%\renewcommand{\arraystretch}{1.4}
%
\begin{tabular}{ cc @{\qquad} l @{\qquad} l}
    \toprule
    \multicolumn{2}{@{}c@{\qquad}}{Output Mode}
  & % blank
  & % blank
    \\
    %---------------------------------------------
    XYZ
  & Axi
  & Variables
  & CGNS Variable Names
    \\
    %---------------------------------------------
    \midrule
    \checkmark
  & \checkmark
  & $\ReyBar{\rho}$,
    $\ReyBar{p}$,
    $\ReyBar{T}$,
    $\ReyAve{\ReyFlx{T}\ReyFlx{T}}$
  & \texttt{Density},
    \texttt{Pressure},
    \texttt{Temperature},
    \texttt{TT}
    \\[0.40em]
    %---------------------------------------------
    \checkmark
  & \checkmark
  & $\ReyBar{T_t}$,
    $\ReyBar{p_t}$
  & \texttt{StagnationTemperature},
    \texttt{StagnationPressure}
    \\[0.40em]
    %---------------------------------------------
    \checkmark
  & \checkmark
  & $\ReyBar{u},\ReyBar{v},\ReyBar{w}$
  & \texttt{VelocityX, VelocityY, VelocityZ}
    \\[0.40em]
    %---------------------------------------------
    % blank
  & \checkmark
  & $\ReyBar{u}_r,\ReyBar{u}_\theta$
  & \texttt{VelocityR, VelocityTheta}
    \\[0.40em]
    %---------------------------------------------
    \checkmark
  & % blank
  & $\ReyAve{\ReyFlx{u  }\ReyFlx{u  }}$,
    $\ReyAve{\ReyFlx{v  }\ReyFlx{v  }}$,
    $\ReyAve{\ReyFlx{w  }\ReyFlx{w  }}$
  & \texttt{uu, vv, ww}
    \\[0.40em]
    %---------------------------------------------
    % blank
  & \checkmark
  & $\ReyAve{\ReyFlx{u  }\ReyFlx{u  }}$,
    $\ReyAve{\ReyFlx{\ur}\ReyFlx{\ur}}$,
    $\ReyAve{\ReyFlx{\ut}\ReyFlx{\ut}}$
  & \texttt{uu, rr, ff}
    \\[0.40em]
    %---------------------------------------------
    \checkmark
  & % blank
  & $\ReyAve{\ReyFlx{u  }\ReyFlx{v  }}$,
    $\ReyAve{\ReyFlx{u  }\ReyFlx{w  }}$,
    $\ReyAve{\ReyFlx{v  }\ReyFlx{w  }}$
  & \texttt{uv, uw, vw}
    \\[0.40em]
    %---------------------------------------------
    % blank
  & \checkmark
  & $\ReyAve{\ReyFlx{u  }\ReyFlx{\ur}}$,
    $\ReyAve{\ReyFlx{u  }\ReyFlx{\ut}}$,
    $\ReyAve{\ReyFlx{\ur}\ReyFlx{\ut}}$
  & \texttt{ur, uf, rf}
    \\[0.40em]
    %---------------------------------------------
    \checkmark
  & % blank
  & $\ReyAve{\ReyFlx{u  }\ReyFlx{T}}$,
    $\ReyAve{\ReyFlx{v  }\ReyFlx{T}}$,
    $\ReyAve{\ReyFlx{w  }\ReyFlx{T}}$
  & \texttt{uT, vT, wT}
    \\[0.40em]
    %---------------------------------------------
    % blank
  & \checkmark
  & $\ReyAve{\ReyFlx{u  }\ReyFlx{T}}$,
    $\ReyAve{\ReyFlx{\ur}\ReyFlx{T}}$,
    $\ReyAve{\ReyFlx{\ut}\ReyFlx{T}}$
  & \texttt{uT, rT, fT}
    \\
    %---------------------------------------------
    \bottomrule
\end{tabular}
%
\endgroup
\end{center}
\end{table}

%%%%%%%%%%%%%%%%%%%%%%%%%%%%%%%%%%%%%%%%%%%%%%%%%%%%%%%%%%%%%%%%%%%%%%%%%%%%%%%%

\SeeAlso{
\inp[nmlTimeAve]{nml\_time\_ave}
}

%%%%%%%%%%%%%%%%%%%%%%%%%%%%%%%%%%%%%%%%%%%%%%%%%%%%%%%%%%%%%%%%%%%%%%%%%%%%%%%%

%\vspace{1.0em}%
%\FIXME{
%\begin{itemize}
%\setlength\topsep{0.0em}
%\setlength\itemsep{0.0em}
%\setlength\parsep{0.0em}
%%
%\item \raggedright
%For rotating reference frame, need more clarity on the reference frame
%and unit vectors of the output variables.
%Are $\fvs{u}=(u,v,w)$ really
%$\mvs{u}=(\mix{u},\mix{v},\mix{w}) = \fms{\FrameRotMatVar}\fvs{u}$?
%Is $\ReyAve{u  }$ really $\ReyAve{\mix{u}        }$?
%Is $\ReyAve{u u}$ really $\ReyAve{\mix{u} \mix{u}}$?
%What are $(\ur,\ut)$?
%%
%\item \raggedright
%Since $\fms{\FrameRotMatVar}=\fms{\FrameRotMatVar}(t)$, the
%rotating frame unit vectors vary with time.
%If $\fms{\FrameRotMatVar}$ is inside the time-average, where do the
%unit vectors of those averaged values point?
%%
%\end{itemize}
%%
%} %end FIXME

%%%%%%%%%%%%%%%%%%%%%%%%%%%%%%%%%%%%%%%%%%%%%%%%%%%%%%%%%%%%%%%%%%%%%%%%%%%%%%%%

}%END_newcommand_scope

